%-------------------------
% Resume in LaTeX
% Anubhav Saini
% Based on Jake Gutierrez's template (https://github.com/jakegut/resume)
% License : MIT
%------------------------

\documentclass[letterpaper,10pt]{article}

\usepackage{latexsym}
\usepackage[empty]{fullpage}
\usepackage{titlesec}
\usepackage{marvosym}
\usepackage[usenames,dvipsnames]{color}
\usepackage{verbatim}
\usepackage{enumitem}
\usepackage[hidelinks]{hyperref}
\usepackage{fancyhdr}
\usepackage[english]{babel}
\usepackage{tabularx}
\usepackage{fontawesome5}

% pdfTeX-only ATS/machine-readability setup (skipped under XeTeX/LuaTeX)
\ifdefined\pdfgentounicode
  \input{glyphtounicode}
\fi

\pagestyle{fancy}
\fancyhf{} % clear all header and footer fields
\fancyfoot{}
\renewcommand{\headrulewidth}{0pt}
\renewcommand{\footrulewidth}{0pt}

% Adjust margins
\addtolength{\oddsidemargin}{-0.6in}
\addtolength{\evensidemargin}{-0.6in}
\addtolength{\textwidth}{1.2in}
\addtolength{\topmargin}{-.7in}
\addtolength{\textheight}{1.4in}

\urlstyle{same}

\raggedbottom
\raggedright
\setlength{\tabcolsep}{0in}

% Sections formatting
\titleformat{\section}{
  \vspace{-4pt}\scshape\raggedright\large
}{}{0em}{}[\color{black}\titlerule \vspace{-5pt}]

% Ensure that generated PDF is machine readable/ATS parsable
\ifdefined\pdfgentounicode
  \pdfgentounicode=1
\fi

%-------------------------
% Custom commands
\newcommand{\resumeItem}[1]{
  \item\small{
    {#1 \vspace{-2pt}}
  }
}

\newcommand{\resumeSubheading}[4]{
  \vspace{-2pt}\item
    \begin{tabular*}{0.97\textwidth}[t]{l@{\extracolsep{\fill}}r}
      \textbf{#1} & \textbf{#2} \\
      \textit{\small#3} & \textit{\small #4} \\
    \end{tabular*}\vspace{-7pt}
}

\newcommand{\resumeProjectHeading}[2]{
    \item
    \begin{tabular*}{0.97\textwidth}{l@{\extracolsep{\fill}}r}
      \small#1 & #2 \\
    \end{tabular*}\vspace{-7pt}
}

\renewcommand\labelitemii{$\vcenter{\hbox{\tiny$\bullet$}}$}

\newcommand{\resumeSubHeadingListStart}{\begin{itemize}[leftmargin=0.15in, label={}]}
\newcommand{\resumeSubHeadingListEnd}{\end{itemize}}
\newcommand{\resumeItemListStart}{\begin{itemize}}
\newcommand{\resumeItemListEnd}{\end{itemize}\vspace{-5pt}}

%-------------------------------------------
%%%%%%  RESUME STARTS HERE  %%%%%%%%%%%%%%%%%%%%%%%%%%%%

\begin{document}

%----------HEADING----------
\begin{center}
    {\Huge \scshape Anubhav Saini} \\ \vspace{1pt}
    \small Vancouver, BC ~ \faPhone\ 604-750-0819 ~
    \href{mailto:sainianubhav01@gmail.com}{\faEnvelope\ \underline{sainianubhav01@gmail.com}} ~
    \href{https://sainianu1.github.io/}{\faGlobe\ \underline{Portfolio}} ~
    \href{https://www.linkedin.com/in/anubhavsainiubc/}{\faLinkedin\ \underline{LinkedIn}} ~
    \href{https://github.com/sainianu1}{\faGithub\ \underline{GitHub}}
    \vspace{-8pt}
\end{center}

%-----------EDUCATION-----------
\section{Education}
  \resumeSubHeadingListStart
    \resumeSubheading
      {Bachelor of Applied Science in Engineering Physics}{Expected May 2027}
      {University of British Columbia -- Specialization in Electrical Engineering and Robotics}{Vancouver, BC}
  \resumeSubHeadingListEnd
  \vspace{-8pt}

%-----------EXPERIENCE-----------
\section{Experience}
  \resumeSubHeadingListStart

    \resumeSubheading
      {Electronics Hardware Intern}{January 2026 -- August 2026}
      {Kardium Inc.}{Burnaby, BC}
      \resumeItemListStart
        \resumeItem{Designed the \textbf{2nd generation of Kardium's highest-volume production PCB} -- a \textbf{6-layer} board carrying both high-voltage and low-voltage signals -- \textbf{increasing its HIPOT withstand strength} for medical HV safety compliance.}
        \resumeItem{Re-defined PCB design rules and fabrication drawings and updated the \textbf{6-layer} design to a new vendor's specifications while maintaining \textbf{IPC Class 3} requirements, \textbf{cutting production costs} on the safety-critical, high-volume board.}
        \resumeItem{Executed layout for Kardium's next generation of \textbf{medical-grade flex PCBs}, routing \textbf{1500V+ signals} within strict safety clearances to deliver short-term cost reduction and rapid drop-in integration into the current production system.}
        \resumeItem{Spearheaded \textbf{end-to-end validation of a 12-layer medical-grade PCB}: authored and executed \textbf{25+ tests across 5 diagnostic suites} (Power, Signal Integrity, Thermal, ADC, Isolation) for \textbf{100\% test coverage}, building custom \textbf{10x/100x coaxial probes} to capture high-fidelity signal-integrity and isolation measurements across the board.}
        \resumeItem{Debugged multiple \textbf{14-layer boards} within a \textbf{10+ board system} via high-speed signal analysis and precision rework.}
      \resumeItemListEnd

    \resumeSubheading
      {Hardware Engineering Intern}{May 2025 -- December 2025}
      {Arlo Technologies}{Vancouver, BC}
      \resumeItemListStart
        \resumeItem{Designed a cost-effective \textbf{Power-ORing Switch} (integrated FW + auxiliary HW control) with \textbf{98\% efficiency} to charge Arlo's camera from the best available power source; applied transistor/diode array logic to create a discrete HV input system.}
        \resumeItem{\textbf{Led R\&D of a new solar-powered security camera with an embedded solar panel}, defined system architecture (configuration of multiple power inputs, finding Boost/charger/fuel gauge ICs), conducted competitive benchmarking, and validated design through comprehensive voltage, power, efficiency, and battery charge/discharge cycle testing.}
        \resumeItem{Investigated and eliminated PIR noise induced by RF radiation from 2.4 GHz Wifi, \textbf{cutting yield loss from 30\% to less than 1\%} and ensuring reliable motion detection during all wireless camera operating states (idle, streaming, recording).}
        \resumeItem{Increased hardware test coverage for PIR and Ambient Light Sensors \textbf{from 60\% to 100\%} to build a coexistence test suite.}
      \resumeItemListEnd

    \resumeSubheading
      {Electrical Engineering Intern}{September 2024 -- December 2024}
      {Sarcomere Dynamics Inc.}{Vancouver, BC}
      \resumeItemListStart
        \resumeItem{Expedited firmware for an in-house fingertip force sensor \textbf{5-10 times cheaper} than on-shelf solutions, reconstructing Normal Force and Shear Direction from raw tri-axis magnetometer (MLX90393) data with \textbf{more than 90\% accuracy} in real time.}
        \resumeItem{Devised a computationally light calibration using pseudo-inverse matrix mapping, verified via K-Means clustering.}
        \resumeItem{Engineered a \textbf{Motor Driver PCB} in \textbf{Altium} around an \textbf{STM32F412} -- SPI to the BLDC motor driver, I2C magnetometer feedback for rotor positioning, CAN for multi-board scalability, and \textbf{3.3V} on-board linear regulation from a single 5V input.}
      \resumeItemListEnd

    \resumeSubheading
      {Product Engineering Intern}{May 2024 -- August 2024}
      {Microchip Technology}{Burnaby, BC}
      \resumeItemListStart
        \resumeItem{Programmed Python test scripts and custom firmware to characterize the SERDES I3C Pad, executing \textbf{8 AC and 5 DC tests} across \textbf{5 voltage/temperature (VT) corners} with per-test BERT/DCSU configuration and pass/fail yield checks.}
        \resumeItem{Formulated and deployed a Multi-threaded testing procedure to \textbf{reduce testing time by more than 40\%} -- inverted the test loop to hold each VT corner fixed while threaded Python maintained uninterrupted voltage and temperature control.}
        \resumeItem{\textbf{Reduced power consumption by 15\%} via systematic voltage/temperature regularization during characterization.}
      \resumeItemListEnd

  \resumeSubHeadingListEnd
  \vspace{-16pt}

%-----------PERSONAL PROJECTS-----------
\section{Personal Projects}
    \resumeSubHeadingListStart

      \resumeProjectHeading
          {\textbf{Solar Express -- DC Power Optimizer} $|$ \emph{Power Electronics, SPICE, Control Loop Design, Python, MPPT}}{}
          \resumeItemListStart
            \resumeItem{Architected a \textbf{600W} panel-level DC-DC power optimizer around a \textbf{4-switch buck-boost} with \textbf{100V GaN FETs at 250kHz}, spanning a 10-80V input / 5-80V output envelope at \textbf{15A}; defined the full spec from PV-panel and string-inverter constraints.}
            \resumeItem{Validated the power stage via automated NGSpice switching simulations across a \textbf{24-point} $V_{in} \times V_{out}$ grid, achieving \textbf{99.5\% pass-through} and \textbf{99.2\% boost efficiency}; verified NEC 690.12 rapid shutdown to \textbf{under 1V in 1.1s} (2s spec).}
            \resumeItem{Performed small-signal plant identification on cycle-averaged SPICE models -- characterized the boost-mode \textbf{RHP zero (3.9kHz)} and built a \textbf{Type-II compensator} with \textbf{$\geq$94$^{\circ}$ phase margin}, \textbf{$\geq$10dB gain margin}, and \textbf{0.12\%} step deviation.}
            \resumeItem{Optimized component selection via parametric loss-budget sweeps ($R_{DS(on)}$ 3-22m$\Omega$, $f_{sw}$ 100kHz-1MHz, L 3.3-22$\mu$H) to lock a \textbf{6.8$\mu$H / 3.2m$\Omega$ / 250kHz} design point; built a \textbf{Python + NGSpice} pipeline validated within \textbf{$\pm$0.5\%} of an analytic model.}
          \resumeItemListEnd

      \resumeProjectHeading
          {\textbf{Automated Robot Race-Car} $|$ \emph{C/C++, Embedded Systems, Analog Circuit Design, PID}}{}
          \resumeItemListStart
            \resumeItem{Built circuitry integrating sensors (IMU, magnetometer, self-made rotary encoder for wheel-spin position mapping) with an STM32, applying RC filtering to cut motor noise and fuse real-time orientation and location data with \textbf{over 85\% accuracy}.}
            \resumeItem{Applied PID control to autonomously track a preset optimal path and developed multi-sensor recalibration routines to re-localize the car after \textbf{1.5-foot ramp jumps} and obstacle encounters, \textbf{achieving the fastest lap time by 4 seconds}.}
          \resumeItemListEnd

    \resumeSubHeadingListEnd
    \vspace{-16pt}

%-----------TECHNICAL SKILLS-----------
\section{Technical Skills}
 \begin{itemize}[leftmargin=0.15in, label={}]
    \small{\item{
     \textbf{PCB Design}: Altium Designer, Schematic Capture, 6-14 Layer \& Flex Layout, High-Voltage Design (HIPOT, IPC Class 3), DFM \\ \vspace{1pt}
     \textbf{Power \& Validation}: Power Electronics (DC-DC), Signal Integrity, SPICE, Oscilloscopes \& Probing, Soldering \& Rework \\ \vspace{1pt}
     \textbf{Embedded \& Software}: STM32, SPI, I2C, UART, CAN, Python, C/C++, MATLAB, TensorFlow, OpenCV, ROS, Linux \\
    }}
 \end{itemize}

\end{document}
